\section{Global inputs and their dependencies}\label{gi:section}

For a standardized real law $P$ with finite $\rho(P)=\E|X|^3$, write
\[
 R_n(P)=\frac{\sqrt n\,\Delta_n(P)}{\rho(P)},\qquad
 C_n=\sup_P R_n(P).
\]
The supremum is over all such laws, without a support restriction.
We use the classical iid moment envelope \cite{Shevtsova2013}
\begin{equation}\label{gi:moment-envelope}
 R_n(P)\le T(\rho):=
 \min\{.469,\ .3322(1+.429/\rho),\ .3031(1+.646/\rho)\}.
\end{equation}
All terminating decimals denote rational numbers.

We include the variational starting point and contact argument;
compare He--Cheng~\cite[Lemmas~6.1 and~6.3, equations~(12)--(15)]{HeCheng2026}.
\begin{lemma}[A contrary extremizer]\label{gi:extremizer}
If $C_n\ge\CE$, the supremum defining $C_n$ is attained by a
standardized law $P$ with $\rho(P)<2$. After reflection if necessary,
there is a finite $H$ such that
\[
 F_n(H)-\Phi(H/\sqrt n)=C_n\rho(P)/\sqrt n,
 \qquad F_n(H)=P^{*n}(( -\infty,H]).
\]
\end{lemma}
\begin{proof}
The classical moment envelope in \eqref{gi:moment-envelope} gives
$C_n\le.469<1/2$ and
$T(2)\le.3031(1+.646/2)=.4010013<\CE$.
Thus a maximizing sequence whose scores tend to $C_n\ge\CE$
eventually has third moments below $2$. Tightness and this third-moment
bound give a weakly convergent subsequence $P_j\Rightarrow P$ whose
first two moments remain $0,1$. Pass also to
$\rho(P_j)\to\rho_*\in[1,2]$, so $\rho(P)\le\rho_*$.

Reflection converts a negative discrepancy into a limit of positive
closed discrepancies. Choose corresponding thresholds $H_j$ with
an error tending to zero. These thresholds stay bounded: Chebyshev's
inequality and the normal tail bound make the discrepancy tend uniformly
to zero as $|H_j|\to\infty$. Passing to $H_j\to H$, the closed-set
inequality for weak convergence gives
\[
 F_n(H)-\Phi(H/\sqrt n)\ge C_n\rho_*/\sqrt n.
\]
Maximality and $\rho(P)\le\rho_*$ force equality and
$\rho(P)=\rho_*$. The moment envelope then gives $\rho(P)<2$.

For the contact condition, put $z=H/\sqrt n$ and $m=\E[X|X|]$.
Perturb the common marginal to
$Q_\varepsilon=(1-\varepsilon)P+\varepsilon\delta_x$.
Write $\mu_\varepsilon,\sigma_\varepsilon^2$ and
$\beta_\varepsilon$ for its mean, variance and centered third
absolute moment. The threshold $H$ stays fixed in these original
coordinates. By the definition of $C_n$,
\[
 J(\varepsilon)=Q_\varepsilon^{*n}(( -\infty,H])
 -\Phi\left(\frac{H-n\mu_\varepsilon}{\sigma_\varepsilon\sqrt n}\right)
 -\frac{C_n\beta_\varepsilon}{\sqrt n\,\sigma_\varepsilon^3}
 \le0=J(0).
\]
The right derivative at zero therefore gives
\begin{align}\label{gi:contact}
 \mathcal D(x)={}&n\{F_{n-1}(H-x)-F_n(H)\}
 +\phi(z)\left\{\sqrt n x+\frac z2(x^2-1)\right\}\notag\\
 &-\frac{C_n}{\sqrt n}
 \left\{|x|^3-\frac32\rho x^2-3mx+\frac12\rho\right\}\le0.
\end{align}
Here $F_0$ is the CDF of the point mass at zero when $n=1$.
Indeed the perturbed mean and variance have derivatives $x$ and
$x^2-1$, and the centered third absolute moment has derivative
$|x|^3-\rho-3mx$. The convolution derivative is exactly
$n\{F_{n-1}(H-x)-F_n(H)\}$, preserving the common marginal.
The integral of $\mathcal D$ under $P$ is zero.
Since $\mathcal D$ is upper semicontinuous, it is zero at every
point of the support, including atoms.

\end{proof}

The heat-smoothing estimate below is already present in
Rodziewicz--Ko\l{}odziejek~\cite[Section~15]{RK2026}.
\begin{lemma}[A coarse Gaussian entry]\label{gi:gaussian-entry}
For $n\ge2$,
\[
 C_n\le C_{n-1}\sqrt{\frac n{n-1}}
       +\frac{\sqrt n}{15(n-1)^{3/2}}
       +\frac{2\sqrt n}{25(n-1)^2}.
\]
In particular, $C_{n-1}\le\CE$ and $n\ge26$ imply
\begin{equation}\label{gi:small-excess}
 n(C_n-\CE)<\frac3{10},\qquad C_n<.422.
\end{equation}
\end{lemma}
\begin{proof}
Put $v=n-1$ and $F_v(t)=\Phi(t/\sqrt v)$. Conditioning on the last
summand, replacing the first $v$ summands by a Gaussian costs at most
$C_v\rho/\sqrt v$. Taylor expansion in the last standardized summand gives
\[
 \left|\E F_v(H-X)-F_v(H)-\tfrac12F_v''(H)\right|
 \le\frac{\phi(0)\rho}{6v^{3/2}}<\frac{\rho}{15v^{3/2}}.
\]
The heat identity $\partial_vF_v=F_v''/2$ gives the matching expansion
in the variance variable, with error at most
$\|\phi'''\|_\infty/(8v^2)<2/(25v^2)$. Here Gaussian Fourier
inversion gives $\|\phi'''\|_\infty\le2/\pi<16/25$.
Combine the errors, multiply by $\sqrt n/\rho$, use $\rho\ge1$,
and take the threshold and law suprema. This is the elementary
heat-smoothing estimate also used in the supporting global route.
For the last assertion put $x=n/(n-1)$. Under the sharp predecessor,
\[
 n(C_n-\CE)\le\frac{\CE x}{\sqrt x+1}
       +\left(\frac1{15}+\frac2{25\sqrt{n-1}}\right)x^{3/2}.
\]
When $n\ge26$, $x\le26/25$, $\sqrt x<51/50$ and $\CE<.41$.
Also $x/(\sqrt x+1)\le(3x+1)/8$, since
$(\sqrt x-1)^2(3\sqrt x+1)\ge0$. Consequently the right side is less than
\[
 \frac{41}{100}\frac{3(26/25)+1}{8}
 +\left(\frac1{15}+\frac2{125}\right)\frac{26}{25}\frac{51}{50}
 =\frac{747107}{2500000}<\frac3{10}.
\]
Also $\CE+3/(10\cdot26)<.422$.
In particular, the sharp bound $C_{29}\le\CE$ already implies
$C_{30}<\CE+.01$, the coarse premise of input (II)
below. The sharp index-$30$ conclusion still needs its own proof.
\end{proof}

The finite-index completion uses the following three conclusions of the
supporting proofs.

\begin{proposition}[Global inputs]\label{gi:inputs}
The low-index companion, actual-index bridge and uniform local comparison
give the following statements. Put
$\delta_n=n\pos{C_{n-1}-C_n}$.
\begin{enumerate}
\item[(I)] $C_j\le\CE$ for $1\le j\le25$;
\item[(II)] if $C_{30}<\CE+1/100$, then $C_j\le\CE$ for
every $31\le j\le46$;
\item[(III)] for every integer $n\ge47$, the conditions
\[
 0\le n(C_n-\CE)\le7/10,\qquad \delta_n=0,
 \qquad C_n<.4247,\qquad C_{n-1},C_{n-2}<.4253
\]
cannot hold simultaneously.
\end{enumerate}
This is the zero-allowance case of the original local exclusion;
the least-failure argument below supplies exactly this condition.
The exclusion (III) is uniform in $n$; its premises will be
discharged by the Gaussian entry and predecessor induction.
\end{proposition}
The precise companion statements and their continuous proof sources
are identified in Appendix~\ref{app:companions}.

% Supporting source map, retained for preparation of the companion package.
% Low indices: work/berry-esseen-n1-base-20261024/n1_base.tex;
% work/berry-esseen-n2-public-check-20261024/collect_replay_20261024.py;
% work/berry-esseen-n3-public-check-20261024/collect_replay_20261024.py;
% work/berry-esseen-n4-analytic-check-20261025/check_analytic_20261025.py;
% work/berry-esseen-n4-public-check-20261024/check_full_closure_exact_20261025.py.
% For the low-index continuous proof map, including the actual LOW,
% conserved-identity, resource and count derivations, see
% work/berry-esseen-c5-complete-20261025/fresh-low-work-20261025/source/
% work/cutoff5_20261020/audit/ANALYTIC_DEPENDENCIES_V2.txt.
% Later low-index interfaces:
% work/berry-esseen-c5-complete-20261025/five_sum_completion_20261025.tex;
% work/berry-esseen-n6-complete-20261025/six_sum_interface_20261025.tex;
% work/berry-esseen-n7-complete-20261025/seven_sum_interface_20261025.tex;
% work/berry-esseen-n8-complete-20261025/eight_sum_interface_20261025.tex;
% work/berry-esseen-n9-complete-20261025/nine_sum_interface_20261025.tex;
% work/berry-esseen-n10-complete-20261025/n10-interface/main.tex;
% work/berry-esseen-n11-14-complete-20261026/conditional_11_14_interface_20261026.tex;
% work/berry-esseen-n15-19-complete-20261026/n15-19-interface/main.tex;
% work/berry-esseen-n20-complete-20261025/twenty_sum_interface_20261025.tex;
% work/berry-esseen-n20-singlebatch-20261026/actual20_singlebatch_interface_20261026.tex;
% work/berry-esseen-n21-24-complete-20261026/n21-24-interface/main.tex;
% work/berry-esseen-n25-full-parent-20261026/n25-interface/main.tex;
% work/berry-esseen-n25-full-parent-20261026/literal25-fourier-proof/main.tex.

The classical inputs are Prawitz's smoothing inequality \cite{Prawitz1972} in the form
given by Shevtsova \cite[Lemma~4.1, p.~263]{Shevtsova2012},
Shevtsova's finite iid moment bounds \cite{Shevtsova2013}, and
Schulz's strict bound for every binomial law \cite[Theorem~1]{Schulz2016}.
In particular, the moment envelope \eqref{gi:moment-envelope} gives
$T(\rho)\le819/2000<\CE$ for $\rho\ge7361/4000$,
excluding the entire larger-moment ray.
Prawitz smoothing is applied at each actual integer, retaining its
endpoint terms and infinite Gaussian tail. The contact, moment-defect,
resource and convolution-count reductions are proved in the companions;
they are not additional unnamed literature assumptions.

Lemma~\ref{gi:extremizer} supplies the common starting point on a
contrary branch, before any use of higher moments. At this same law,
\begin{equation}\label{gi:predecessor}
 n\pos{R_{n-1}(P)-C_n}\le n\pos{C_{n-1}-C_n}.
\end{equation}
Thus a global predecessor bound licenses contact allowances throughout
the argument. A score bound restricted to a moment band is never
substituted for that predecessor.

For (I), the first four direct bounds are
\[
 C_1\le\tfrac25,\quad C_2\le\tfrac{2\sqrt2}{7},\quad
 C_3\le\tfrac{59\sqrt3}{250},\quad C_4\le.4097<\CE.
\]
The companion compositions use the unrestricted predecessor in order
through $24$. At $25$, a support-free Fourier bound excludes
$\rho\ge6/5$; the two lower bands use the global $C_{24}<.419$
and their complete resource/count covers. Their continuous reductions,
unseeded resource steps and closed arithmetic covers are the precise
inputs listed in Appendix~\ref{app:companions}. The fifth-index proof
is also given separately in \cite{RKFive2026}; subsequent source
developments are collected in \cite{RKSources2026}.

% Actual-index bridge sources:
% work/berry-esseen-n31-primary-check-20261025/sharp31_bridge_composition_20261025/main.tex;
% work/berry-esseen-n35-complete-20261025/actual35_conditional_composition_20261025.tex;
% work/berry-esseen-n36-38-complete-20261026/actual36_38_conditional_composition_20261026.tex;
% work/berry-esseen-n39-complete-20261026/actual39_two_step_composition_20261026/main.tex;
% work/berry-esseen-n40-route-20261026/n40-48-interface/main.tex.

For (II), the coarse bound at $30$ gives allowance $31/100$ at
a contrary index-$31$ extremizer. The eight closed moment bands have
actual-$31$ support, resource and count proofs. The bridges through
$34$ keep $mD_m\le31d$, acquire support at each actual $m$, and use
that integer's binomial law. Subsequent implications proceed through
$35$--$39$ and $40$--$46$, retaining both deleted predecessors where
required. Thus the bridge starts from the stated coarse premise
without any monotonicity assumption on $C_n$.

% Large-index source:
% work/berry-esseen-n49-check-20261025/n49_completion_20261025.tex;
% substantive saved-output joins:
% work/berry-esseen-n49-check-20261025/join_n49_fresh_results_20261025.py;
% work/berry-esseen-n49-check-20261025/check_n49_final_join_20261025.py.

For (III), the moment envelope gives $\rho<7361/4000$ at a contrary
extremizer. Corollary~\ref{gc:tail-support} then gives radius $3.84$
directly from contact and the deleted-predecessor allowance.
It also supplies the two left radial contractions and the central
radius-$2.5$ contraction, in their original acquisition order.
Thus the six initial index-band and subsequent radial support
calculations are replaced by the same cubic bound. The bounded-threshold
coordinate-strip exclusions retain their original proofs.
The resource and confinement schedule then supplies $45$ central
and $38$ left probability intervals, covering $[.304,.47]$ after
all complementary intervals have been excluded. Their resources
retain the actual weighted quantities $nD$ and $n\E U^2$.
The complete local comparison is uniform for every $n\ge47$.
At positive noise it gives a strict loss relative to this law's
own Bernoulli score; at zero noise Schulz's strict binomial theorem
applies directly. The acquisition includes both deleted-score caps
and the full excess interval $0\le n(C_n-\CE)\le.7$.

\begin{proposition}[Induction from two sharp predecessors]\label{gi:tail}
Input~(III) and $C_{45},C_{46}\le\CE$ imply
$C_n<\CE$ for every integer $n\ge47$.
\end{proposition}
\begin{proof}
If a failure exists, choose the least $n\ge47$ with $C_n\ge\CE$.
Both predecessors satisfy $C_{n-1},C_{n-2}\le\CE$, so $\delta_n=0$
and both deleted-score caps in (III) hold. The Gaussian entry gives
\[
 0\le n(C_n-\CE)<3/10<7/10,\qquad
 C_n<.422<.4247.
\]
All premises of (III) are now satisfied, a contradiction.
\end{proof}

The moment-defect method and the original finite subdivisions retain
their attribution to Rodziewicz--Ko\l{}odziejek \cite{RK2026}.
The conditional continuation in Section~11 of that working manuscript
retains analytic and historical certificate premises, including the
sharp assertions at $25$--$30$. It does not by itself establish (I) or
fill the intervening finite range. Inputs (I)--(III) here refer to the
companion proofs with their continuous reductions, complete exact
calculations and closed-domain compositions. In particular, the proof
order is (I), the new actual-index steps $26$--$30$, (II) through
$46$, and Proposition~\ref{gi:tail}. The five finite ceilings in
the independent $n\ge49$ route are unnecessary for this composition.
